\section{Support, margin, and growth}
\label{sec:introduction}

A negative direction can exist and yet be difficult to resolve.
For a reflected exponential form, a pair away from the imaginary
axis creates an indefinite contribution. Other locations can mask
that contribution on a short interval, or leave a negative value
whose size is much smaller than the norm of its test. This leads
to three different questions: when does negativity first occur;
how much support is needed for a specified normalized margin;
and how rapidly can the negative envelope grow for one fixed
infinite configuration?

The distinction is already visible in a single pair. If
\(H_f(z)=\int f(u)e^{zu}\,du\), the pair at \(\beta,-\beta\)
contributes \(2\Re(H_f(\beta)\overline{H_f(-\beta)})\).
On \((-L,L)\), its largest normalized negative value is
\[
 \beta^{-1}\{\sinh(2\beta L)-2\beta L\}.
\]
It is positive for every \(L>0\), but for \(\beta L\ll1\)
it is approximately \((4/3)\beta^2L^3\). Thus existence of a
negative test and existence of a unit-margin test are already
different assertions before a background is added.

We study these questions for locally finite divisors \(Z\) in
\(|\Re z|\le a\), invariant under \(z^\#=-\overline z\), with
positive integer multiplicities and a polynomial bound on the
number of locations counted with multiplicity. The form is
\[
 Q_Z(f,g)=\sum_{z\in Z}\nu(z)H_f(z)\overline{H_g(z^\#)}.
\]
We work throughout with actual compact smooth tests. Write
\(\cM_Z(L)\) for the supremum of the positive part of
\(-Q_Z(f,f)/\|f\|_2^2\) on \((-L,L)\), and \(r_Z(\eta)\)
for the infimum of radii admitting a test with margin at least
\(\eta\). The definitions, including infinite envelopes and
nonattained infima, are given in Section~\ref{sec:basics}.

\subsection{The uniform support law}
Our first main result fixes the counting parameters and varies
the displacement and requested margin. Suppose
\[
 N_0(Y)=\sum_{|\Im z|\le Y}\nu(z)\le C(1+Y)^m
 \quad(Y\ge0),\qquad C\ge2,
\]
and require \(\beta,-\beta\in Z\). All other reflected pairs
are permitted. Let \(\cR_\eta(C,\beta;a,m)\) be the worst
value of \(r_Z(\eta)\) over this entire class. For fixed
\(a,C,m\), we prove
\begin{equation}\label{intro:frontier}
 \cR_\eta(C,\beta;a,m)\asymp_{a,C,m}
 \beta^{-1}\log\left(1+(\beta\eta)^{1/(2\lfloor C\rfloor-1)}\right),
 \qquad \eta\ge1,\quad0<\beta\le\min(a,1).
\end{equation}
This is an order estimate, not an exact leading minimax constant.
At unit margin the support exponent is
\((2\lfloor C\rfloor-2)/(2\lfloor C\rfloor-1)\).
The witness may depend on the entire divisor. The conclusion
does not provide a predetermined finite collection of tests.

The local integer budget determines the exponent, while the
rest of the counting bound enters the constants. Arbitrarily
close distinct nodes prevent a naive interpolation estimate
based on inverse separation. We instead prescribe opposite
values on two reflected groups of nearby nodes, annihilate
an intermediate band, and suppress the full remaining tail.
The grouped contour length saves one power in the coefficient
estimate. Applying the intermediate polynomial as a differential
operator in physical coordinates, before a probability
convolution, preserves the needed norm bound. Finite unit-weight
screens give the matching obstruction against every test.

An additional conjugation symmetry about the target ordinate
changes the exponent denominator to \(4\lfloor C/2\rfloor-1\).
This is a parity effect of the integer count: an origin atom
is invisible to an odd test and the remaining local mass is
even. It is not legitimate to infer the centered symmetry
from global conjugation after translating a nonzero target
height. Section~\ref{sec:frontier} proves both versions with
their full parameter ranges.

\subsection{A finite transition and two infinite-divisor conclusions}
Finite screens make the mechanism more explicit. For the
Chebyshev screen
\[
 \mathcal S_{n,\beta}
 =128\sum_{j=1}^n
       \left[i\beta n\cos\frac{(2j-1)\pi}{2n}\right]
       +[\beta]+[-\beta],
 \qquad d=2\lfloor n/2\rfloor+1,
\]
we compute the unique negative eigenvalue and obtain
\[
 \cM_{\mathcal S_{n,\beta}}(L)
 =c_n\beta^{2d}L^{2d+1}\bigl(1+O_n((\beta L)^2)\bigr),
 \qquad c_n>0.
\]
Here \(n\) is fixed and \(\beta L\to0\). For any fixed
polynomial-count axis divisor \(A\), the leading transition
survives:
\begin{equation}\label{intro:fixed-background}
 \cM_{\mathcal S_{n,\beta}+A}(x\beta^{-q})
 \longrightarrow c_nx^{2d+1},\qquad
 q=\frac{2d}{2d+1},\quad\beta\downarrow0,
\end{equation}
locally uniformly for \(x>0\). The background may define
an unbounded positive form. It is held fixed, and the
window expands; neither qualification can be omitted.
The proof in Section~\ref{sec:screening} uses finite smooth
moment recovery followed by a separate high-frequency
estimate on that fixed space.

For a fixed divisor with a nonempty off-axis part, put
\(s_Z=\sup\{\Re z:z\in Z,\ \Re z>0\}\). Polynomial counting
alone gives
\begin{equation}\label{intro:lower-growth}
 \liminf_{L\to\infty}\frac{\log\cM_Z(L)}L\ge2s_Z.
\end{equation}
The logarithm is interpreted in the extended sense if an
envelope is infinite. If, in addition,
\[
 B_Z=\sup_{y\in\R}
 \sum_{\substack{z=\sigma+i\gamma\in Z\\
                 \sigma>0,\ |\gamma-y|\le1}}
       \nu(z)\sigma^2<\infty,
\]
then \(\log\cM_Z(L)/L\to2s_Z\), and
\(r_Z(\eta)\sim(2s_Z)^{-1}\log\eta\) as \(\eta\to\infty\).
Section~\ref{sec:growth} proves the lower construction in
every sufficiently large window and the upper estimate by
a weighted area-mean argument. The condition on \(B_Z\)
is sufficient, not a characterization of local semiboundedness.

Our second main conclusion shows how far local semiboundedness
alone is from a growth law. Given any finite-valued positive
nondecreasing \(G\) on \([2,\infty)\), we construct one
divisor, invariant under both reflection and conjugation,
in \(|\Re z|\le1\), for which
\begin{equation}\label{intro:arbitrary}
 \begin{gathered}
 N_0(Y)\le Y+2,\qquad N_t(1)\le2+\log(2+|t|),\\
 \cM_Z(L)<\infty\quad(0<L<\infty),\qquad
 \cM_Z(L)\ge G(L)\quad(L\ge2).
 \end{gathered}
\end{equation}
The constants in the count bounds do not depend on \(G\).
The construction separates two radii with a finite screen
and then places conjugate pairs of screens far apart in
height. It controls both the effect of each new component
on old witnesses and the effect of old components on the
new witness. Summable small-window negative bounds coexist
with the prescribed large-window witnesses.

There is no conflict between \eqref{intro:frontier},
\eqref{intro:fixed-background}, and \eqref{intro:arbitrary}.
The first takes a supremum over varying divisors with fixed
counting parameters; the second holds one axis background
fixed while a displacement tends to zero; the third fixes
a specially constructed infinite divisor and lets support
grow. Keeping these quantifiers visible is central to the
subject.

\subsection{Relation to prior work and the contribution here}
The arithmetic motivation is Weil's positivity criterion,
but the forms in this paper are abstract. Bombieri
\cite[Sections 4, 8--11]{Bombieri} studies compact-support
variational problems for the actual Weil functional and
finite exponential-kernel matrices. In particular, his
finite-matrix analysis determines negative inertia from
off-axis pairs and discusses the delicate passage to
infinite configurations. We use the elementary finite
inertia mechanism as background, not as a new result.
Our uniform result specifies the joint support and
\(L^2\)-margin cost over a polynomial-count class, allowing
an arbitrary additional off-axis part and controlling its
entire tail. It is not an arithmetic strengthening of
Bombieri's conclusions.

The collision powers in finite exponential systems also
have substantial precedents in superresolution. Demanet
and Nguyen \cite[Theorem 2 and Corollary 3]{DemanetNguyen}
obtain sharp small-bandwidth powers for sparse recovery
on a grid. Batenkov, Demanet, Goldman, and Yomdin
\cite[Theorems 3.2 and 3.9]{BatenkovEtAl} give lower bounds
for nonuniform Fourier matrices under specified clustering
conditions, together with examples having the matching
cluster-size power. Those are
singular-value and recovery problems. Here the quantity
is an indefinite, normalized quadratic value, and the
upper construction must preserve signs in an infinite
complex reflected configuration. The appearance of an
odd power is not itself a novelty.

Likewise, the moment and Schur-complement calculations
used for a finite screen belong to established
perturbation theory. Barthelm\'e and Usevich
\cite{BarthelmeUsevich} treat kernel matrices in the flat
limit, and Usevich and Barthelm\'e \cite{UsevichBarthelme}
treat analytic symmetric perturbations, including
indefinite matrices. The contribution of our finite-screen
analysis is the explicit negative-margin transition and
the fixed infinite-background limit
\eqref{intro:fixed-background}, not a new general
matrix-perturbation method. We include the calculation
needed here rather than require a general perturbation
theorem. Classical interpolation and orthogonal-polynomial
normalizations are recorded in \cite{DLMF}.

Finite zero-annihilating products multiplied by high powers
of sinc also occur in Zhu's RH-conditional construction
\cite[Section 13.2]{Zhu}. That product--window template is
not claimed as novel in this work. The sign-compatible interpolation,
physical norm estimate, and full-tail bounds required for
our indefinite problem are proved below; none of that
preprint's other arithmetic assertions is used.

The principal conclusions established here are therefore the
full-class fixed-count law \eqref{intro:frontier}, its
centered-symmetry refinement, the stated fixed-background
transition, and the separation between local
semiboundedness and arbitrary growth
\eqref{intro:arbitrary}. The fixed-divisor growth estimates
connect these conclusions and give a useful sufficient
class in which the outer displacement is recoverable
from the exponential rate.

Recent work on compact Weil operators and near-radical
vectors, for example Connes and Consani
\cite{ConnesConsani}, concerns additional arithmetic
structure that our abstract divisor class does not
possess. The final section gives the actual prime,
gamma, and polar expression in our normalization and
explains what an arithmetic upper bound
would imply. We prove no such new upper bound here.